\subsection{Matlis duality}

In this section, we assume that the ring A is local Noetherian.

DEFINITION.--- An A-module I is said to be a Matlis A-module if it is injective, if \(\mathfrak{m}_{\mathrm{A}}\) is its unique associated prime ideal and if the \(\kappa_{\mathrm{A}}\) -vector space \(\operatorname{Hom}_{\mathrm{A}}(\kappa_{\mathrm{A}}, \mathrm{I})\) has dimension 1.

Let \(e: \kappa_{\mathrm{A}} \to I\) be an injective envelope of \(\kappa_{\mathrm{A}}\) (A, X, p. 20, th. 2). The A-module I is a Matlis module, and every Matlis A-module is isomorphic to I (no. 2, cor. of prop. 3). If A is a discrete valuation ring with field of fractions K, the A-module K/A is a Matlis module (no. 1, example). If A is a local Artinian ring, the A-module A is a Matlis module if and only if A is a Gorenstein ring (§ 3, no. 7, lemma 1).

Let I be a Matlis A-module. For every integer \(n \geqslant 0\) let us denote by \(I_n\) the sub-A-module of I consisting of the elements annihilated by \(\mathfrak{m}_A^n\) . By prop. 2 of no. 2, the A-module I is the union of the \(I_n\), the A-module \(I_1\) is of length 1 (that is, isomorphic to \(\kappa_A\) ) and the A-module I is an injective envelope of \(I_1\); moreover, the canonical homomorphism of graded gr(A)-modules

\[
\beta : \mathrm{gr} ^ {\mathrm{m} _ {\mathrm{A}}} (\mathrm{I}) \longrightarrow \operatorname{Homgr} _ {\kappa_ {\mathrm{A}}} (\mathrm{gr} (\mathrm{A}), \mathrm{I} _ {1})\tag{3}
\]

is an isomorphism. By prop. 3 of no. 2, the A-module structure of I extends uniquely to a structure of \(\widehat{\mathrm{A}}\) -module, and the canonical homomorphism \(\widehat{\mathrm{A}} \to \operatorname{End}_{\mathrm{A}}(\mathrm{I})\) is bijective.

Lemma 3.-- Let I be a Matlis A-module. Then:

\begin{enumerate}
\def\labelenumi{\alph{enumi})}
\item
  I is a A-hat-Matlis module;
\item
  the A-module I is Artinian and a cogenerator (A, X, p. 18, def. 3).
\end{enumerate}

Since the A-module I is injective, the \(\mathbf{A} / \mathfrak{m}_{\mathbf{A}}^{n}\) -module \(\mathrm{I}_n\) is injective for each \(n\) ( \(n^{\circ}2\), (lemma 1, a)). Since \(\mathrm{I}_n\) is the set of elements of I annihilated by \(\mathfrak{m}_{\widehat{\mathrm{A}}}^{n}\) , the \(\widehat{\mathrm{A}}\) -module I is injective (lemma 1, b)). It is indecomposable over \(\widehat{\mathrm{A}}\) since it is so over \(\mathrm{A}\); as it contains the sub- \(\widehat{\mathrm{A}}\) -module \(\mathrm{I}_1\) isomorphic to \(\kappa_{\mathrm{A}}\) , one has \(\mathfrak{m}_{\widehat{\mathrm{A}}} \in \operatorname{Ass}_{\widehat{\mathrm{A}}}(\mathrm{I})\) , hence \(\operatorname{Ass}_{\widehat{\mathrm{A}}}(\mathrm{I}) = \{\mathfrak{m}_{\widehat{\mathrm{A}}}\}\) (prop. 1), whence a).

Let us now prove that I is Artinian. To every sub-A-module M of I, associate the graded ideal \(\mathfrak{a}_{\mathrm{M}}\) of gr(A) defined as follows: an element of gr(A) \(_n\) belongs to \((\mathfrak{a}_{\mathrm{M}})_{n}\) if it is annihilated by all the linear forms \(\boldsymbol{\beta}(x)\), where x ranges over \(\left((\mathrm{M}\cap\mathrm{I}_{n+1})+\mathrm{I}_{n}\right)/\mathrm{I}_{n}\). Let M and N be submodules of I such that \(N\subset M\); one has \(a_{M}\subset a_{N}\) . Suppose \(a_{M}=a_{N}\) ; one has \((M\cap I_{n+1})+I_{n}=(N\cap I_{n+1})+I_{n}\) for all n since \(\beta\) is an isomorphism. By induction on n we deduce \(M\cap I_{n+1}=N\cap I_{n+1}\) for all n, whence finally M=N.

That being so, let \(M_{0} \supset M_{1} \supset \ldots \supset M_{n} \supset \ldots\) be a decreasing sequence of sub-A-modules of I; the increasing sequence \(a_{M_{0}} \subset a_{M_{1}} \subset \ldots\) is stationary, since \(\text{gr}(A)\) is an \(\kappa_{A}\) -algebra of finite type. The sequence \((M_{i})_{i \geqslant 0}\) is therefore stationary, which implies that the A-module I is Artinian. Finally, the A-module I is a cogenerator by virtue of A, X, p. 18, prop. 12.

Let M be an A-module. Recall (A, VIII, § 4, n° 6) that the socle of M is the sum of the simple submodules of M, that is, the set of elements of M annihilated by \(\mathfrak{m}_{\mathrm{A}}\); it is a \(\kappa_{\mathrm{A}}\) -vector space, canonically isomorphic to \(\operatorname{Hom}_{\mathrm{A}}(\kappa_{\mathrm{A}}, \mathrm{M})\).

Lemma 4. Let I be a Matlis module over A and let M be an \(A\)-module. The following conditions are equivalent:

\begin{enumerate}
\def\labelenumi{(\roman{enumi})}
\item
  M is Artinian;
\item
  every element of M is annihilated by a power of \(\mathfrak{m}_{\mathrm{A}}\), and the socle of M has finite dimension over \(\kappa_{\mathrm{A}}\);
\item
  there exists an integer \(n \geqslant 0\) and an injective A-linear map from M into \(\mathrm{I}^n\) .
\end{enumerate}

When these conditions are satisfied, every injective envelope of M is isomorphic to \(I^{s}\), where s is the dimension over \(\kappa_{A}\) of the socle of M.

(iii)⇒(i): this is clear since the A-module I is Artinian (Lemma 3).

(i)⇒(ii): suppose M is Artinian. Let \(x \in M\); the decreasing sequence of submodules \(m_{A}^{n}x\) of M is stationary. Let n be an integer such that \(m_{A}^{n+1}x = m_{A}^{n}x\); Nakayama's lemma implies \(m_{A}^{n}x = 0\). Moreover, the socle of M is Artinian as an A-module, hence also as a \(\kappa_{A}\)-vector space, which means it is of finite dimension.

(ii)⇒(iii): suppose condition (ii) holds; let \(e:M\to J\) be an injective envelope of M. We have \(\operatorname{Ass}(M)\subset\{\mathfrak{m}_{A}\}\), hence \(\operatorname{Ass}(J)\subset\{\mathfrak{m}_{A}\}\) ( \(n^{\circ} 1\), see Remark 2), and J is isomorphic to I \(^{(c)}\) for a cardinal c ( \(n^{\circ} 1\), th. 1). Let \(x\) be a nonzero element of J annihilated by \(\mathfrak{m}_{A}\); as the A-module Ax is simple and its intersection with \(e(M)\) is not reduced to 0, \(x\) belongs to \(e(M)\). Thus \(e\) induces an isomorphism from the socle of M onto that of J; consequently, the socle of M has dimension c, which proves (iii) as well as the last assertion.

Lemma 5.--- Every Artinian A-hat-module is Artinian as a Λ-module.

Let M be a \(\widehat{\mathbf{A}}\) -module is Artinian; every element of M is annihilated by a power of \(\mathfrak{m}_{\widehat{\Lambda}}\), hence by a power of \(\mathfrak{m}_{\mathbf{A}}\). By lemma 2 of no. 2, the sub-A-modules of M are its sub- \(\widehat{\mathbf{A}}\) -modules, so M is Artinian as an A-module.

Now fix a Matlis A-module I. For every A-module M, denote by \(D_{\mathrm{A}}(\mathrm{M})\) the \(\widehat{A}\) -module

\[
\mathrm{D} _ {\mathrm{A}} (\mathrm{M}) = \operatorname{Hom} _ {\mathrm{A}} (\mathrm{M}, \mathrm{I}).
\]

The \(\widehat{\mathrm{A}}\) -module \(D_{\mathrm{A}}(A)\) is canonically identified with I, the \(\widehat{\mathrm{A}}\) -module \(D_{\mathrm{A}}(\mathrm{I})\) with \(\widehat{\mathrm{A}}\) ( \(n^{\circ}\) 2, prop. 3), and the \(\widehat{\mathrm{A}}\) -module \(D_{\mathrm{A}}(\kappa_{\mathrm{A}})\) with \(I_{1}\) (loc. cit.).

For any A-linear map \(f: M \rightarrow N\), we shall denote

\[
\mathrm{D} _ {\mathrm{A}} (f): \mathrm{D} _ {\mathrm{A}} (\mathrm{N}) \rightarrow \mathrm{D} _ {\mathrm{A}} (\mathrm{M})
\]

the map \(\widehat{\mathrm{A}}\) -linear \(\operatorname{Hom}_{\Lambda}(f,1_{\mathrm{I}})\) . Since the \(\Lambda\) -module I is injective, the sequence \((\mathrm{D}_{\Lambda}(g),\mathrm{D}_{\mathrm{A}}(f))\) is exact for every exact sequence \((f,g)\) of A-linear maps.

\((D_{\hat{A}}(y), D_{\hat{A}}(f))\) is exact for every exact sequence \((f, y)\) of maps.

We will apply these definitions to the ring \(\hat{A}\) endowed with the Matlis module I (lemma 3, a)); for every \(\hat{A}\) -module P, \(D_{\hat{A}}(P)\) is therefore the sub- \(\hat{A}\) -module \(\text{Hom}_{\hat{A}}(P, I)\) of \(D_{A}(P)\). It amounts to the same thing to say that P is artinian as an A-module or as \(\hat{A}\) -module (Lemma 5); if this is the case, we have \(D_{\hat{A}}(P) = D_{A}(P)\) (loc. cit.).

Let M be an A-module. For \(m \in M\), the map \(f \mapsto f(m)\) of \(\mathrm{D}_{\mathrm{A}}(\mathrm{M})\) in I is \(\widehat{\mathrm{A}}\) -linear; denote it by \(\alpha_{\mathrm{M}}(m)\). One thus defines an A-linear homomorphism

\[
\alpha_ {\mathrm{M}}: \mathrm{M} \longrightarrow \mathrm{D} _ {\widehat {\mathrm{A}}} (\mathrm{D} _ {\mathrm{A}} (\mathrm{M})) .
\]

One denotes by \(\widehat{\alpha}_{\mathrm{M}}:\widehat{\mathrm{A}}\otimes_{\mathrm{A}}\mathrm{M}\longrightarrow \mathrm{D}_{\widehat{\mathrm{A}}}(\mathrm{D}_{\mathrm{A}}(\mathrm{M}))\) the \(\widehat{\mathrm{A}}\) -linear map deduced from \(\alpha_{\mathrm{M}}\)

THEOREM 2.--- Let M be an A-module.

\begin{enumerate}
\def\labelenumi{\alph{enumi})}
\item
  For M to be Artinian, it is necessary and sufficient that the \(\hat{\mathrm{A}}\)-module \(\mathrm{D}_{\mathrm{A}}(\mathrm{M})\) be of finite type. When this is the case, the homomorphism \(\alpha_{\mathrm{M}}\) is bijective.
\item
  For M to be of finite type, it is necessary and sufficient that \(D_{A}(M)\) be Artinian (as an A-module or as \(\hat{A}\)-module). In this case the homomorphism \(\hat{\alpha}_{M}\) is an isomorphism.
\item
  For M to be of finite length, it is necessary and sufficient that \(\mathrm{D}_{\mathrm{A}}(\mathrm{M})\) be of finite length (as an A-module or as \(\widehat{A}\)-module). In this case \(\alpha_{M}\) is an isomorphism from M onto \(\mathrm{D}_{\mathrm{A}}(\mathrm{D}_{\mathrm{A}}(\mathrm{M}))\) and one has \(\mathrm{long}_{\mathrm{A}}(\mathrm{D}_{\mathrm{A}}(\mathrm{M})) = \mathrm{long}_{\mathrm{A}}(\mathrm{M})\).
\end{enumerate}

Let us first prove that the homomorphism \(\alpha_{\mathrm{M}}\) is injective for every A-module M. Let \(m\) be a nonzero element of M; its annihilator is contained in \(\mathfrak{m}_{\mathrm{A}}\). There therefore exists a surjective A-homomorphism from \(Am\) onto \(\kappa_{\mathrm{A}}\) , and consequently a nonzero homomorphism from \(Am\) into I. Since I is injective, it extends to a homomorphism \(f:M\to I\) such that \(f(m)\neq0\). This proves the injectivity of \(\alpha_{\mathrm{M}}\).

Suppose the A-module M is Artinian. By Lemma 4, there exists an integer \(r\) and an injective A-linear map \(f: \mathrm{M} \to \mathrm{I}^r\) The homomorphism \(\mathrm{D}_{\mathrm{A}}(f): \mathrm{D}_{\mathrm{A}}(\mathrm{I}^r) \longrightarrow \mathrm{D}_{\Lambda}(\mathrm{M})\) is then surjective; since \(\mathrm{D}_{\mathrm{A}}(\mathrm{I}^r)\) is identified with \(\widehat{\mathrm{A}}^r\) , this proves that the \(\widehat{\mathrm{A}}\) -module \(\mathrm{D}_{\mathrm{A}}(\mathrm{M})\) is of finite type. Analogously, if M is of finite type, there exists an integer \(n\) and a surjective homomorphism \(u: \mathrm{A}^n \to \mathrm{M}\) ; the homomorphism \(\mathrm{D}_{\Lambda}(u): \mathrm{D}_{\mathrm{A}}(\mathrm{M}) \to \mathrm{I}^n\) is injective, so that \(\mathrm{D}_{\mathrm{A}}(\mathrm{M})\) is artinian (as an A-module or as a \(\widehat{\mathrm{A}}\) -module).

Suppose now that the \(\widehat{\mathrm{A}}\) -module \(\mathrm{D}_{\mathrm{A}}(\mathrm{M})\) is Artinian; the same holds for the \(\widehat{\mathrm{A}}\) -module \(\mathrm{D}_{\widehat{\mathrm{A}}}(\widehat{\mathrm{A}} \otimes_{\mathrm{A}} \mathrm{M})\) which is canonically isomorphic to it. By the preceding, the \(\widehat{\mathrm{A}}\) -module \(\mathrm{D}_{\widehat{\mathrm{A}}}(\mathrm{D}_{\widehat{\mathrm{A}}}(\widehat{\mathrm{A}} \otimes_{\mathrm{A}} \mathrm{M}))\) is of finite type, and the same holds for \(\widehat{\mathrm{A}} \otimes_{\mathrm{A}} \mathrm{M}\) which is isomorphic to a submodule of \(\mathrm{D}_{\widehat{\mathrm{A}}}(\mathrm{D}_{\widehat{\mathrm{A}}}(\widehat{\mathrm{A}} \otimes_{\mathrm{A}} \mathrm{M}))\)). Consequently M is a finitely generated A-module (I, § 3, no. 6, prop. 11 and III, § 3, no. 5, prop. 9). Likewise, if \(\mathrm{D}_{\mathrm{A}}(\mathrm{M})\) is a \(\widehat{\mathrm{A}}\) -module of finite type, \(\mathrm{D}_{\widehat{\mathrm{A}}}(\mathrm{D}_{\mathrm{A}}(\mathrm{M}))\) is a \(\widehat{\mathrm{A}}\) By the preceding, it is an Artinian module, hence an Artinian A-module (Lemma 5), and the same holds for M. Finally, the modules of finite length are the finitely generated Artinian modules (A, VIII, § 1, no. 1, Prop. 1), hence \(\mathrm{D}_{\mathrm{A}}(\mathrm{M})\) has finite length if and only if M has finite length.

Suppose M is Artinian. There exists an integer r and an injective A-linear map \(f : M \to I^{r}\); since I is Artinian (Lemma 3), the A-module Coker(f) is also Artinian, and one can find an integer s and an exact sequence of A-modules

\[
0 \to \mathrm{M} \stackrel {f} {\longrightarrow} \mathrm{I} ^ {r} \stackrel {g} {\longrightarrow} \mathrm{I} ^ {s}.
\]

We deduce a commutative diagram with exact rows

\[
\begin{array}{c c c c c c} 0 \longrightarrow & \mathrm{M} & \xrightarrow {f} & \mathrm{I} ^ {r} & \xrightarrow {g} & \mathrm{I} ^ {s} \\ & \alpha_ {\mathrm{M}} \Biggl \downarrow & & \alpha_ {\mathrm{I} ^ {r}} \Biggl \downarrow & & \alpha_ {\mathrm{I} ^ {s}} \Biggl \downarrow \\ 0 \to & \mathrm{D} _ {\widehat {\mathrm{A}}} (\mathrm{D} _ {\mathrm{A}} (\mathrm{M})) & \xrightarrow {\mathrm{D} _ {\widehat {\mathrm{A}}} (\mathrm{D} _ {\Lambda} (f))} & \mathrm{D} _ {\widehat {\mathrm{A}}} (\mathrm{D} _ {\Lambda} (\mathrm{I} ^ {r})) & \xrightarrow {\mathrm{D} _ {\widehat {\mathrm{A}}} (\mathrm{D} _ {\Lambda} (g))} & \mathrm{D} _ {\widehat {\Lambda}} (\mathrm{D} _ {\Lambda} (\mathrm{I} ^ {s})) \end{array} .
\]

The \(\widehat{\mathbf{A}}\)-module \(\mathrm{D}_{\widehat{\mathbf{A}}}(\mathrm{D}_{\mathbf{A}}(\mathrm{I}))\) is identified with I and \(\alpha_{\mathbf{I}}\) with the identity map; consequently \(\alpha_{\mathbf{I}^r}\) and \(\alpha_{\mathbf{I}^s}\) are bijective, and the same holds for \(\alpha_{\mathbf{M}}\) (A, X, p. 7, cor. 3).

If the A-module M is finitely generated, there exist integers m and n and an exact sequence of A-modules

\[
\mathrm{A} ^ {m} \longrightarrow \mathrm{A} ^ {n} \longrightarrow \mathrm{M} \rightarrow 0;
\]

We deduce a commutative diagram with exact rows:

\[
\begin{array}{c c c c c} \widehat {\mathrm{A}} ^ {m} & \longrightarrow & \widehat {\mathrm{A}} ^ {n} & \longrightarrow & \widehat {\mathrm{A}} \otimes_ {\mathrm{A}} \mathrm{M} \longrightarrow 0 \\ \widehat {\alpha} _ {\mathrm{A} ^ {m}} \Bigg \downarrow & & \widehat {\alpha} _ {\mathrm{A} ^ {n}} \Bigg \downarrow & & \widehat {\alpha} _ {\mathrm{M}} \Bigg \downarrow \\ \widehat {\mathrm{A}} ^ {m} & \longrightarrow & \widehat {\mathrm{A}} ^ {n} & \longrightarrow & \mathrm{D} _ {\mathrm{A}} (\mathrm{D} _ {\mathrm{A}} (\mathrm{M})) \to 0. \end{array}
\]

As \(\widehat{\alpha}_{\mathrm{A}}\) is equal to \(1_{\widehat{\mathrm{A}}}\), it follows that \(\widehat{\alpha}_{\mathrm{M}}\) is an isomorphism.

It remains to prove the equality. \(\operatorname{long}_{\mathrm{A}}(\mathrm{M}) = \operatorname{long}_{\mathrm{A}}(\mathrm{D}_{\Lambda}(\mathrm{M}))\) when M is of finite length. We may assume \(M \neq 0\) ; there then exists an exact sequence

\[
0 \rightarrow \kappa_ {\mathrm{A}} \longrightarrow \mathrm{M} \longrightarrow \mathrm{N} \rightarrow 0,
\]

from which one deduces an exact sequence

\[
0 \rightarrow \mathrm{D} _ {\mathrm{A}} (\mathrm{N}) \longrightarrow \mathrm{D} _ {\mathrm{A}} (\mathrm{M}) \longrightarrow \mathrm{D} _ {\mathrm{A}} (\kappa_ {\mathrm{A}}) \rightarrow 0.
\]

We have long \(_{A}\)(M) = long \(_{A}\)(N) + 1 and

\[
\operatorname{long} _ {\Lambda} \left(\mathrm{D} _ {\Lambda} (\mathrm{M})\right) = \operatorname{long} _ {\mathrm{A}} \left(\mathrm{D} _ {\mathrm{A}} (\mathrm{N})\right) + \operatorname{long} _ {\mathrm{A}} \left(\mathrm{D} _ {\mathrm{A}} \left(\kappa_ {\mathrm{A}}\right)\right) = \operatorname{long} _ {\mathrm{A}} \left(\mathrm{D} _ {\mathrm{A}} (\mathrm{N})\right) + 1;
\]

we conclude by induction on the integer \(_{A}\) (M).

Remark.--- Suppose the ring A is Artinian. We have \(\text{long}_{\text{A}}(\text{I}) = \text{long}_{\text{A}}(\text{D}_{\text{A}}(A)) = \text{long}(\text{A})\) (th. 2, c)). Let M be a finitely generated A-module; it admits an injective envelope isomorphic to I \(^{s}\) , where s is the dimension of the socle of M (lemma 4). Consequently we have \(\text{long}_{\text{A}}(\text{M}) \leqslant s \text{ long}(\text{A})\); for equality to hold, it is necessary and sufficient that M be injective. In particular, for the A-module A to be injective, it is necessary and sufficient that its socle have dimension 1; we thus recover lemma 1 of § 3, no. 7.

